\documentclass{article}
\usepackage[a4paper, total={6in, 8in}]{geometry}
\usepackage{amssymb, amsmath, parskip}
 % Required for inserting images

\title{CENG60002: Reaction Engineering II}
\begin{document}
\maketitle
\section{Topic 0. Conditions for Phase Stability, and Reaction Equilibrium}

Chemical reactors are not at equilibrium, as the net rates of reactions are usually not allowed to become zero. Considering a single-phase multicomponent system, if we declare the following variables: \(T\), temperature, \(P\), pressure, and \(n_i\) the number of moles of each component \(i\). The Gibbs energy of the system G utilizes these state variables to declare the conditions for equilibrium, 
\begin{align*}
    dG = -SdT + VdP + \sum_{j} \mu_j dn_j
\end{align*}
where \(S\) is the entropy, \(V\) is the volume; both of the system, and \(\mu\) is the chemical potential of component \(j\).



\section{Topic 1. External Mass Transfer}

The analysis of external mass transfer will be done on a catalyst surface. We define the following variables: \(C_{A,s}\), the concentration of reactant A in the catalyst, \(C_{A,ss}\): the concentration of reactant A at the surface of the catalyst \(z=0\), and finally \(C_A\): the concentration of A in the bulk liquid. 

We can define the following rates:
\begin{enumerate}
    \item \textbf{Mass Transfer}: \(r_{A,mass} = k_{m_c}\left(C_A - C_A^{surf}\right)\), where \(k_{m_c} = \frac{D_{AB}}{\delta}\) (derived from Fick's law of diffusion), this assumes a linear diffusion profile along the laminar film across the surface of the catalyst.
    \item \textbf{Reaction Rate}: \(r_{A,kinetics} = k_{surf} C_A^{surf}\); simple rate of reaction, where \(k_{surf} = k_0 \exp^{\left(-\frac{E_a}{RT}\right)}\)
\end{enumerate}

At equilibrium:
\begin{align*}
r_{A,\text{mass}} &= r_{A,\text{kinetics}} \\
k_{m_c}\left(C_A - C_A^{\text{surf}}\right) &= k_{\text{surf}} C_A^{\text{surf}} \\
C_{A}^{surf} &= \left(\frac{k_{m,c}}{k_{m,c} + k_{surf}}\right) C_A
\end{align*}

Finally, we define the overall rate constant \(k_0\) as
\begin{align*}
\frac{1}{k_0} &= \frac{1}{k_{m_c}} + \frac{1}{k_{surf}}
\end{align*}

Due to this construction, we can observe the following limiting cases:
\begin{enumerate}
    \item \text{Rapid Mass Transfer}: \(k_{m_c} >> k_s\), resulting in \(k_0 \approx k_s\), and \(C_{A,surf} \approx C_A\), which is therefore rate-controlled.
    \item \text{Rapid Reaction}: \(k_{s} >> k_{m_c}\), resulting in \(k_0 \approx k_{m_c}\), and \(C_{A,surf} \approx 0\), which is therefore diffusion-controlled.
\end{enumerate}

There are also several important dimensionless numbers:
\begin{enumerate}
    \item \textbf{Fick's Law Corrections}: \(\delta_A = \frac{\left(r + s + ...\right) - \left(a + b + ...\right)}{a}\).
    \item \textbf{Reynolds Number}: \(Re = \frac{G d_p}{\mu} = \frac{inertial\;effects}{viscous\;effects}\).
    \item \textbf{Prandlt Number}: \(Re = \frac{C_p \mu}{\lambda} = \frac{viscous\;effects}{heat\;transfer\;effects}\).
    \item \textbf{Schmidt Number}: \(Sc = \frac{\mu}{\rho_f D} = \frac{viscous\;effects}{diffusion\;effects}\).
\end{enumerate}

\end{document}
